% =====================================================================
%  Fluffy (Google CTF 2025) - reversing
%  Engenharia reversa de um app Flutter/Dart AOT e otimizacao de
%  forca bruta via periodicidade de 16 rodadas
%
%  Compilacao:  pdflatex slides.tex   (rode 2x para o sumario/overlays)
%  Necessita:   beamer, tikz, listings (padrao em TeX Live / MiKTeX)
% =====================================================================
\documentclass[aspectratio=169]{beamer}

% - idioma e codificacao -------------------------------------------
\usepackage[T1]{fontenc}
\usepackage[utf8]{inputenc}
\usepackage[brazilian]{babel}

% - matematica / extras --------------------------------------------
\usepackage{amsmath,amssymb}
\usepackage{graphicx}
\usepackage{xcolor}
\usepackage{listings}

% - tikz -----------------------------------------------------------
\usepackage{tikz}
\usetikzlibrary{arrows.meta, calc, positioning}

% - tema -----------------------------------------------------------
\usetheme{Madrid}
\usecolortheme{whale}
\setbeamertemplate{navigation symbols}{}

% - identidade da apresentacao --------------------------------------
\title[Fluffy - Google CTF 2025]{Quebrando um cofre em Flutter: engenharia reversa e otimização de força bruta}
\subtitle{Desafio \emph{Fluffy} (Google CTF 2025) --- Dart AOT, token derivado de timestamp e periodicidade de 16 rodadas}
\author[Integrantes do grupo]{{\footnotesize
  \begin{tabular}{@{}r@{\hspace{0.8em}}l@{}}
    822070 & Nicolas Queiroz Bertozzo\\
    821402 & Gustavo Cesar Bento Laurindo\\
    811537 & Danilo da Silva Pinto\\
    811943 & Lucas de Oliveira Rodrigues Alves\\
    813583 & Lucas de Araújo Cardoso\\
    823127 & Nicole Correa Ramos\\
    811649 & Gabriel Vianna de Oliveira Spolon\\
    812110 & Pedro Guilherme Torres das Neves\\
    832627 & Lucas Martinez Domingues
  \end{tabular}}}
\institute[UFSCar]{Universidade Federal de São Carlos}
\date{2 de outubro de 2026}

% - configuracoes de listagem --------------------------------------
\lstset{
  language=Python,
  basicstyle=\ttfamily\footnotesize,
  keywordstyle=\color{blue!60!black}\bfseries,
  commentstyle=\color{gray!60!black}\itshape,
  stringstyle=\color{teal},
  showstringspaces=false,
  columns=fullflexible,
  keepspaces=true,
  breaklines=true,
  frame=single,
  backgroundcolor=\color{gray!8},
}

% =====================================================================
\begin{document}

% ---------------------------------------------------------------------
\begin{frame}[plain]
  \begin{beamercolorbox}[sep=12pt,center,wd=\textwidth]{title}
    \usebeamerfont{title}\inserttitle\par
    \vspace{0.25em}
    {\usebeamerfont{subtitle}\insertsubtitle\par}
  \end{beamercolorbox}
  \vspace{0.7em}
  \begin{center}
    {\large\usebeamercolor[fg]{structure}\textbf{Integrantes do grupo}\par}
    \vspace{0.4em}
    {\footnotesize
    \begin{tabular}{@{}r@{\hspace{1.2em}}l@{}}
      822070 & Nicolas Queiroz Bertozzo\\
      821402 & Gustavo Cesar Bento Laurindo\\
      811537 & Danilo da Silva Pinto\\
      811943 & Lucas de Oliveira Rodrigues Alves\\
      813583 & Lucas de Araújo Cardoso\\
      823127 & Nicole Correa Ramos\\
      811649 & Gabriel Vianna de Oliveira Spolon\\
      812110 & Pedro Guilherme Torres das Neves\\
      832627 & Lucas Martinez Domingues
    \end{tabular}\par}
    \vspace{0.7em}
    {\normalsize\textbf{\insertinstitute}\par}
    \vspace{0.35em}
    {\small \insertdate\par}
  \end{center}
\end{frame}

% ---------------------------------------------------------------------
\begin{frame}{Roteiro}
  \begin{enumerate}
    \item O desafio e a motivação: um cofre que só encripta
    \item Fundamentos: onde mora o código em um app Flutter (Dart AOT)
    \item Reconstrução: âncoras de string e o algoritmo (token + cifra)
    \item As fragilidades: token determinístico e período de 16 rodadas
    \item O ataque otimizado e o custo computacional
    \item Resultados: os três segredos e a flag
    \item Lições aprendidas e referências
  \end{enumerate}
\end{frame}

% ---------------------------------------------------------------------
\begin{frame}{O desafio e a motivação}
  \begin{columns}[T]
    \begin{column}{0.52\textwidth}
      \textbf{Por que este desafio?}
      \begin{itemize}
        \item App \emph{Flutter} que \alert{só encripta} (decriptar está ``em
              construção''), com algoritmo próprio (\emph{security by obscurity}).
        \item Encripta com um \texttt{token} e um \texttt{PIN} de 4 dígitos.
      \end{itemize}
      \textbf{Os três segredos (Base62):}
      {\scriptsize
      \begin{block}{}
        \texttt{fmMf7mIMbHcPoQmLGx1CO0XVGBmhjTaYhB0}\\
        \texttt{5O6WRgCajs3QSTyohnu2hldds18mjkx}\\
        \texttt{fgv99dOvazsvEESh7DPKbb3k0I3RW}
      \end{block}}
    \end{column}
    \begin{column}{0.44\textwidth}
      \textbf{Nosso objetivo}
      \begin{itemize}
        \item \alert{Localizar e reconstruir} o algoritmo dentro do binário
              compilado (\texttt{libapp.so}, Dart AOT).
        \item Recuperar os três segredos \emph{sem} conhecer o PIN nem o token.
        \item Demonstrar a \alert{redução da busca} (periodicidade + horários),
              e não apenas força bruta ``com muitos núcleos''.
      \end{itemize}
      \begin{block}{Pista}
        1º segredo exibe os \textbf{segundos}; ``amigo da Suíça''
        $\Rightarrow$ fuso \texttt{Europe/Zurich}.
      \end{block}
    \end{column}
  \end{columns}
\end{frame}

% ---------------------------------------------------------------------
\begin{frame}{Fundamentos: de Dart a \texttt{libapp.so}}
  \begin{itemize}
    \item Flutter \emph{debug} usa \textbf{JIT} (interpretado, \emph{hot
          reload}); Flutter \emph{release} usa \textbf{AOT}: o código Dart é
          compilado para \alert{código de máquina nativo} e embarcado num
          objeto ELF --- o \texttt{libapp.so}.
    \item O \texttt{libapp.so} contém: (1) o código de todas as funções;
          (2) a \emph{snapshot} (\texttt{\_kDartIsolateSnapshotData}) com o
          \emph{object pool}: strings, constantes e metadados; (3) tabelas do GC.
    \item \textbf{Por que é difícil de reverter:}
      \begin{itemize}
        \item \alert{Sem tabela de símbolos} das funções do app;
        \item convenção de chamada e \emph{garbage collector} próprios do Dart;
        \item formato da \emph{snapshot} \alert{versionado} (Dart 3.8.1, snapshot
              \texttt{830f4f59\ldots}).
      \end{itemize}
    \item \textbf{Mas:} strings e constantes ficam \alert{legíveis} na seção de
          dados --- é o ponto de entrada da análise.
  \end{itemize}
\end{frame}

% ---------------------------------------------------------------------
\begin{frame}{Reconstrução: âncoras de string e o algoritmo}
  \begin{columns}[T]
    \begin{column}{0.48\textwidth}
      \textbf{Âncoras (via \texttt{strings})}
      \begin{tabular}{@{}l l@{}}
        \texttt{gctf25\_} & seed do token\\
        \texttt{0123\ldots xyz} & alfabeto Base62\\
        \texttt{Invalid character} & dica no decode\\
      \end{tabular}
      \vspace{0.4em}
      \begin{itemize}
        \item Seguindo as \emph{cross-references} (Ghidra/radare2) localizam-se
              \texttt{rol8}/\texttt{ror8} e o laço de \texttt{pin} rodadas.
        \item \texttt{blutter}/\texttt{darter} automatizam, desde que
              compatíveis com o Dart 3.8.1.
      \end{itemize}
      \begin{block}{Validação (round-trip)}
        Encriptar com \texttt{token}=\texttt{I6X6vyQzRuH},
        \texttt{pin}=8126 reproduz \alert{byte a byte} o ciphertext do 1º
        segredo $\Rightarrow$ reconstrução correta.
      \end{block}
    \end{column}
    \begin{column}{0.52\textwidth}
      \textbf{Token} (derivado do tempo):
      \begin{center}
        {\footnotesize \texttt{token = Base62( SHA1("gctf25\_" + ts)[:8] )}}
      \end{center}
      \textbf{Cifra} (rodada $i$, byte $j$; $s$=segredo, $t$=token):
      {\small
      \[
        s_j = \operatorname{rol8}\!\big((s_j + t_{j \bmod 8}) \bmod 256,\; j \bmod 8\big)
      \]}
      depois: rota o array de $s$ (direita) e de $t$ (esquerda), e
      {\small
      \[
        t_j = \operatorname{ror8}\!\big(t_j,\; (pin \oplus ((i \bmod 4)+1)) \bmod 8\big)
      \]}
    \end{column}
  \end{columns}
\end{frame}

% ---------------------------------------------------------------------
\begin{frame}{O token que não é aleatório}
  \begin{center}
    {\small \texttt{token = Base62( SHA1("gctf25\_" + unix\_seconds)[:8] )}}
  \end{center}
  \begin{itemize}
    \item É uma \alert{KDF determinística do timestamp}: o ``aleatório'' é, na
          verdade, o relógio. Há só $86400$ tokens por dia.
    \item O 1º segredo mostra \texttt{04/08/2023 13:37} \alert{com segundos};
          os demais mostram só o minuto $\Rightarrow$ \textbf{60} candidatos.
    \item A dica ``Suíça'' fixa o fuso \texttt{Europe/Zurich} para reconstruir
          o timestamp (a data exibida é local; o app usa o epoch UTC).
  \end{itemize}
  \begin{center}
    \begin{tabular}{@{}l l l@{}}
      minuto exibido & $\longrightarrow$ & 60 tokens candidatos \\
      fuso Zurique   & $\longrightarrow$ & converte data local em epoch \\
    \end{tabular}
  \end{center}
\end{frame}

% ---------------------------------------------------------------------
\begin{frame}[fragile]{A cifra: rodadas e rotações}
  Essência do \texttt{encrypt} (Python equivalente ao Dart reconstruído):
  \begin{lstlisting}[basicstyle=\ttfamily\scriptsize]
def encrypt(pin, token, secret):          # `pin` rodadas no total
    tok = list(base62_decode(token))
    enc = list(secret)
    for i in range(pin):
        enc = [rol8((s + tok[j % 8]) % 256, j % 8)   # soma + rotacao
               for j, s in enumerate(enc)]
        enc = [enc[-1]] + enc[:-1]          # array do segredo: roda -> direita
        tok = tok[1:] + [tok[0]]            # array do token:  roda -> esquerda
        tok = [ror8(d, (pin ^ ((i & 3) + 1)) % 8) for d in tok]
    return base62_encode(bytes(enc))
  \end{lstlisting}
  \begin{itemize}
    \item O número de rodadas é o \alert{próprio PIN} ($1$ a $9999$) --- por
          isso a força bruta ingênua é cara.
    \item Todas as operações são bijetivas $\Rightarrow$ existe decriptação
          (inverso em ordem reversa).
    \item A rotação de bits \alert{só depende de $pin \bmod 8$ e de $i \bmod 4$}.
  \end{itemize}
\end{frame}

% ---------------------------------------------------------------------
\begin{frame}{O período de 16 rodadas}
  Com $k = pin \bmod 8$, a rotação de bits é
  $r_i = (k \oplus ((i \bmod 4)+1)) \bmod 8$ --- \alert{período 4 em $i$}.
  \begin{columns}[T]
    \begin{column}{0.55\textwidth}
      \begin{center}
      \begin{tikzpicture}[font=\scriptsize, x=0.52cm, y=0.72cm, scale=0.85]
        \foreach \i in {0,...,15}{%
          \pgfmathsetmacro{\rr}{int(mod(\i,4)+1)}%
          \pgfmathsetmacro{\yy}{0.55*\rr}%
          \draw[blue!60!black,thick] (\i,0) -- (\i,\yy);%
          \fill[blue!60!black] (\i,\yy) circle (1.1pt);%
          \node[below=1pt] at (\i,0) {\i};%
          \node[above=1pt,blue!60!black] at (\i,\yy) {\rr};%
        }
        \draw[->] (-0.6,0) -- (16.2,0);
        \node[align=left,anchor=north west,black!65] at (0.4,3.4)
          {$k=0$: sequência $1,2,3,4$ repetida 4 vezes};
        \node[align=left,anchor=north west,blue!60!black] at (0.4,2.7)
          {soma em 16 $=4\times10=40 \equiv 0 \pmod 8$};
        \node[align=left,anchor=north west,red!80!black] at (0.4,2.0)
          {soma em 8 $=2\times10=20 \equiv 4 \pmod 8$};
      \end{tikzpicture}
      \end{center}
    \end{column}
    \begin{column}{0.45\textwidth}
      \begin{tabular}{@{}l c c@{}}
        $k = pin \bmod 8$ & soma em 16 & soma em 8\\
        \hline
        $0,1,2,3$ & $40 \equiv 0$ & $20 \equiv 4$\\
        $4,5,6,7$ & $72 \equiv 0$ & $36 \equiv 4$\\
      \end{tabular}
      \vspace{0.2em}
      \begin{itemize}
        \item Estado do token = (posição do array, rotação acumulada $S_i$).
        \item Repete quando $i\equiv0 \pmod 8$ \textbf{e} $S_i\equiv0 \pmod 8$.
        \item Em $i=8$: $S_8\equiv4$ (\alert{não} repete); em $i=16$:
              $S_{16}\equiv0$ (\alert{repete}).
      \end{itemize}
      \begin{block}{Conclusão}
        O token dinâmico tem \alert{período exato 16} --- só 16 estados por agenda.
      \end{block}
    \end{column}
  \end{columns}
\end{frame}

% ---------------------------------------------------------------------
\begin{frame}{Conjuntos de PIN mod 8}
  Como $r_i$ só depende de $k = pin \bmod 8$, há apenas \alert{8 agendas} de
  rotação. Gerando os 16 estados de cada agenda e agrupando os idênticos:
  \begin{itemize}
    \item \textbf{40 estados} distintos de token no total;
    \item \textbf{8 estados} alcançáveis sob \emph{todos} os 8 valores de $k$;
    \item \textbf{32 estados} alcançáveis sob \emph{exatamente 2} valores de $k$.
  \end{itemize}
  \begin{center}
    \begin{tabular}{@{}l c c@{}}
      & estados & pares (estado, agenda)\\
      \hline
      universais & $8$ & $8 \times 8 = 64$\\
      restritos  & $32$ & $32 \times 2 = 64$\\
      \hline
      total & $40$ & $\mathbf{128}$ \quad (não $40\times8=320$)
    \end{tabular}
  \end{center}
  \begin{block}{Por que importa}
    Ao decriptar de trás para frente não sabemos o estado final do token nem o
    $pin \bmod 8$ --- mas há só \alert{128} combinações relevantes a testar.
  \end{block}
\end{frame}

% ---------------------------------------------------------------------
\begin{frame}[fragile]{O ataque otimizado: caminhar de trás para frente}
  \begin{columns}[T]
    \begin{column}{0.58\textwidth}
      \begin{lstlisting}[basicstyle=\ttfamily\scriptsize]
for tok_pat, pin_mods in tok_patterns:   # <= 40 estados finais
    for pin_mod in pin_mods:             # agendas (8 ou 2)
        dyn, secret = tok_pat[:], list(base62_decode(encr))
        for step in range(max_pin):      # volta 1 rodada/vez
            secret = secret[1:] + [secret[0]]
            secret = [(ror8(e, j % 8) - dyn[j % 8]) % 256
                      for j, e in enumerate(secret)]
            dyn = [rol8(d, INV_ROT[pin_mod][step % 4])
                   for d in dyn]
            dyn = [dyn[-1]] + dyn[:-1]
            if is_printable(secret):     # oraculo: ASCII 0x20-0x7E
                return bytes(secret)
      \end{lstlisting}
    \end{column}
    \begin{column}{0.42\textwidth}
      \begin{itemize}
        \item Pré-computa os 40 estados finais (período 16) e, para cada par
              $(estado, agenda)$, \alert{desfaz} uma rodada por passo.
        \item A tabela \texttt{INV\_ROT} inverte a rotação de bits em $O(1)$.
        \item \textbf{Oráculo}: o primeiro estado todo imprimível
              (\texttt{0x20--0x7E}) é o texto claro; o nº de passos é o PIN.
      \end{itemize}
      \begin{alertblock}{}
        Custo por token: $\mathbf{128 \times 9999 \approx 1{,}28M}$ passos,
        contra $\sum_{p=1}^{9999} p \approx 50M$ da ingênua.
      \end{alertblock}
    \end{column}
  \end{columns}
\end{frame}

% ---------------------------------------------------------------------
\begin{frame}{Custo computacional comparado}
  \begin{center}
  \begin{tabular}{@{}l r r c@{}}
    Estratégia & por token & total ($\times60$) & fator\\
    \hline
    Ingênua (recalcula $p$ rodadas por PIN) & $\approx 50M$ & $\approx 3{,}0\times10^{9}$ & $1\times$\\
    Período 16 + PIN mod 8 & $\approx 1{,}28M$ & $\approx 7{,}7\times10^{7}$ & \alert{$\approx 39\times$}\\
  \end{tabular}
  \end{center}
  \vspace{0.5em}
  \begin{itemize}
    \item A redução \alert{não} muda o nº de combinações de chave
          ($60\times9999$); elimina o \alert{trabalho redundante} de recomputar
          $p$ rodadas a cada PIN.
    \item É exatamente a pergunta da flag: \emph{otimizou, ou só usou muitos
          núcleos?} --- a resposta técnica é a primeira.
  \end{itemize}
  \begin{center}
    {\small Medido (Ryzen 7 5700X3D, 3 processos): \textbf{378\,s}, \textbf{394\,s},
    \textbf{535\,s} por segredo.}
  \end{center}
\end{frame}

% ---------------------------------------------------------------------
\begin{frame}{Resultados: os três segredos e a flag}
  \begin{center}
  \begin{tabular}{@{}l l c l l@{}}
    Data & segundo & PIN & token & texto claro\\
    \hline
    \texttt{04/08/2023 13:37} & 27 & 8126 & \texttt{I6X6vyQzRuH} & \texttt{CTF\{Ok4y\_h4v3\_u\_0ptim1zed\_}\\
    \texttt{07/09/2024 18:52} & 31 & 5178 & \texttt{BMiZFI8Xr3Q} & \texttt{brUt3\_f0rcE\_0R\_y0u\_jUst}\\
    \texttt{03/03/2025 22:07} & 48 & 7490 & \texttt{KZ95PpF1RFq} & \texttt{\_uSeD\_a\_l0t\_0f\_c0Res?\}}\\
  \end{tabular}
  \end{center}
  \vspace{0.4em}
  \begin{block}{Flag completa}
    \texttt{CTF\{Ok4y\_h4v3\_u\_0ptim1zed\_brUt3\_f0rcE\_0R\_y0u\_jUst\_uSeD\_a\_l0t\_0f\_c0Res?\}}
  \end{block}
  \begin{itemize}
    \item O 3º segredo é o mais caro: o segundo correto (48) aparece tarde na
          varredura dos 60.
  \end{itemize}
\end{frame}

% ---------------------------------------------------------------------
\begin{frame}[fragile]{Execução}
  \begin{columns}[T]
    \begin{column}{0.56\textwidth}
      \textbf{Demo} --- \texttt{python fluffy.py demo}:
      \begin{lstlisting}[basicstyle=\ttfamily\scriptsize]
$ python fluffy.py demo
token      : I6X6vyQzRuH
PIN        : 8126
ciphertext : fmMf7mIMbHcPoQmLGx1CO0XVGBmhjTaYhB0
round-trip : b'CTF{Ok4y_h4v3_u_0ptim1zed_'
round-trip OK
ciphertext reproduzido: OK
      \end{lstlisting}
    \end{column}
    \begin{column}{0.44\textwidth}
      \textbf{Ataque completo} (60 s varridos):
      \begin{lstlisting}[basicstyle=\ttfamily\scriptsize]
segredo 1 : 378 s
segredo 2 : 394 s
segredo 3 : 535 s
      \end{lstlisting}
      \vspace{0.4em}
      {\scriptsize O crack otimizado, com o token correto em mãos, recupera o
      texto do 1º segredo em \textbf{6,5 s}. O resto do tempo é a varredura dos
      60 segundos.}
    \end{column}
  \end{columns}
\end{frame}

% ---------------------------------------------------------------------
\begin{frame}[fragile]{Como reproduzir}
  \begin{columns}[T]
    \begin{column}{0.6\textwidth}
      \begin{lstlisting}[basicstyle=\ttfamily\footnotesize]
$ python fluffy.py demo
$ python fluffy.py token '04/08/2023 13:37' 27
      # -> I6X6vyQzRuH
$ python fluffy.py crack \
    '04/08/2023 13:37' \
    'fmMf7mIMbHcPoQmLGx1CO0XVGBmhjTaYhB0'
      # ataque otimizado (~6 min)
      \end{lstlisting}
    \end{column}
    \begin{column}{0.4\textwidth}
      \begin{itemize}
        \item \texttt{fluffy.py} é autocontido (só biblioteca padrão): token,
              Base62, cifra e o crack otimizado.
        \item \texttt{crack.py} é o solver oficial (mesma lógica).
        \item O relatório completo está em \texttt{relatorio-engenharia-reversa.md}.
      \end{itemize}
    \end{column}
  \end{columns}
\end{frame}

% ---------------------------------------------------------------------
\begin{frame}{Lições aprendidas}
  \begin{itemize}
    \item \textbf{Security by obscurity não é segurança}: num binário sem
          símbolos, as strings entregaram o algoritmo inteiro.
    \item \textbf{Chave derivada do tempo é chave fraca}: o ``token aleatório''
          era o relógio --- entropia real (64 bits sorteados) é essencial.
    \item \textbf{Estrutura periódica vaza}: o período de 16 rodadas permitiu
          reaproveitar o ciclo e cortar o custo em $\approx39\times$ --- projetar
          cifra exige pensar na repetição de estado.
    \item \textbf{RE e criptoanálise juntas}: localizar o algoritmo (RE) foi
          metade; entender a matemática da cifra (criptoanálise) foi a outra.
  \end{itemize}
\end{frame}

% ---------------------------------------------------------------------
\begin{frame}{Referências}
  \begin{itemize}
    \item Google CTF 2025 --- \emph{Fluffy} (repositório, fonte e write-up):
          \url{https://github.com/google/google-ctf/tree/1655538e8c8b41451d39f670ef15a5af22979ca9/2025/quals/rev-fluffy}
    \item \texttt{blutter} (reversão de Flutter/Dart AOT):
          \url{https://github.com/worawit/blutter}
    \item \texttt{darter} (parser de snapshot Dart):
          \url{https://github.com/mildsunrise/darter}
    \item \texttt{frida} (instrumentação dinâmica): \url{https://frida.re}
    \item Documentação Dart AOT (\texttt{gen\_snapshot}): \url{https://dart.dev}
  \end{itemize}
  \vspace{0.5em}
  \begin{center}
    {\large Obrigado!}\\[0.2em]
    {\small Perguntas?}
  \end{center}
\end{frame}

\end{document}
